% ==================== 第 6 节：实证与规范（85–102 分钟）====================
\section{实证分析与规范分析：两种不同的问题}

经济学研究常被混作一谈的两类问题，其实属于不同的逻辑类别：\emph{实证}（positive）研究“是什么”，\emph{规范}（normative）研究“应当怎样”。把二者分开，是博士生必须具备的方法论基本功。

\subsection{实证分析：回答“是什么”}

\begin{definition}[实证命题]
\emph{实证命题}是关于世界\emph{是}什么样、以及“若……则……”会怎样变化的可检验陈述。它回答三类问题：是什么（描述）、为什么（机制）、会怎样（预测）。
\end{definition}

实证命题的三个特征：
\begin{enumerate}
  \item \textbf{概念清晰}：每个术语都有可操作的定义（如“正常品”定义为 $\partial x_l^{*}/\partial m>0$）；
  \item \textbf{机制明确}：因果链条在模型内部可追踪（$\theta\to B(\theta)\to C(\theta)$）；
  \item \textbf{预测可检验}：原则上存在数据可以\emph{推翻}它。
\end{enumerate}

\begin{example}[比较静态是典型实证命题]
“收入增加会提高对正常品的需求”——在保持偏好、价格不变且满足假设 $A$ 时，参数 $m$ 的上升导致选择集合发生特定变化：
\[
m'\ge m \;\Longrightarrow\; x_l^{*}(p,m')\ge x_l^{*}(p,m)\quad(\text{正常品}).
\]
又如：“税率 $\tau$ 上升会降低劳动供给”——这是关于 $C(\theta)$ 对 $\tau$ 的比较静态，原则上可用微观数据检验。
\end{example}

实证结论的规范写法是带\emph{条件}的：
\[
\text{在保持 }Z\text{ 不变且满足假设 }A\text{ 时，参数 }p\text{ 的变化导致选择发生某种变化。}
\]
丢掉条件，实证命题就从“可检验”退化为“口号”。

\subsection{规范分析：回答“应当怎样”}

\begin{definition}[规范命题]
\emph{规范命题}是关于“什么应当发生”“哪个结果更好”的陈述。它必须援引某个\emph{价值标准}（福利准则），而价值标准不能从事实中自动推出。
\end{definition}

最常用的价值标准有：
\begin{itemize}
  \item \textbf{帕累托标准}：状态 $x$ 帕累托优于 $y$，若无人变差且至少一人变好（严格写法见下）。帕累托改进得到的是\emph{效率}判断，不涉及分配；
  \item \textbf{功利主义社会福利函数}：
        \[
        W(u_1,\dots,u_n)=\sum_{i=1}^{n}\alpha_i u_i,\qquad \alpha_i\ge 0,
        \]
        把社会评价写成个体效用的加权和；
  \item \textbf{罗尔斯（max–min）标准}：$W=\min_i u_i$，社会评价取决于最不利者。
\end{itemize}

\begin{remark}[权重不是事实]
功利主义中的权重 $\alpha_i$ \emph{不是}由个体选择事实自动推出的，它体现的是评价者的价值立场（谁更重要）。把某个 $\alpha$ 当作“客观的”，正是把规范伪装成实证。同理，效率提高不意味公平改善；观察到某政策“有效”，也不等于该政策“正当”。
\end{remark}

\begin{definition}[帕累托最优]
配置 $x$ 称为\emph{帕累托最优}（Pareto optimal），若不存在其他可行配置 $x'$ 使所有人 $i$ 满足 $u_i(x'_i)\ge u_i(x_i)$ 且至少一个 $j$ 满足 $u_j(x'_j)>u_j(x_j)$。
\end{definition}

\begin{example}[效率与公平可以背离]
把所有资源给一个人而其余人为零，可以是帕累托最优（无法在不损害该人的前提下改善他人），但几乎无人认为它“应当”如此。帕累托标准对分配沉默，这正是它作为\emph{最低限度}价值标准的原因与局限。
\end{example}

\subsection{三层分离}

政策评价应分三层进行：

\begin{figure}[H]
\centering
\begin{tikzpicture}[
  box/.style={draw,thick,rounded corners=2pt,minimum width=3.3cm,minimum height=1.05cm,align=center,fill=gray!8},
  lab/.style={font=\small,align=left,text width=4.6cm}
]
  \node[box] (a) at (0,0) {\textbf{事实层}\\ \small 政策改变了什么？};
  \node[box] (b) at (4.6,0) {\textbf{机制层}\\ \small 经何种约束/激励/均衡传导？};
  \node[box] (c) at (9.2,0) {\textbf{价值层}\\ \small 依何种福利标准评价？};
  \draw[->,very thick] (a) -- (b);
  \draw[->,very thick] (b) -- (c);
  \node[lab] at (0,-1.35) {\small 例：最低工资使低技能就业下降 2\%};
  \node[lab] at (4.6,-1.35) {\small 例：劳动需求弹性与覆盖范围决定效应大小};
  \node[lab] at (9.2,-1.35) {\small 例：用 $W=\sum\alpha_i u_i$ 还是罗尔斯标准？};
\end{tikzpicture}
\caption{政策评价的三层分离}
\end{figure}
\figexplan{三个方框依次是政策评价的三层：事实层（政策改变了什么）、机制层（经何种约束、激励或均衡传导）、价值层（依何种福利标准评价）。箭头表示论证推进方向，每层下方给出最低工资的例子。关键纪律：不得把价值层伪装成事实层——由“提高了效率”推不出“应当实行”。}

\begin{remark}
把第三层（价值）伪装成第一层（事实），是经济学论证中最常见的越界。例如“最低工资提高了效率，所以应该实行最低工资”——前半句即使有证据，后半句仍需额外的价值前提才能成立。
\end{remark}

\subsection{案例演练：把一个政策问题拆成正反两半}

\begin{center}
\small
\begin{tabular}{p{2.6cm}p{5.6cm}p{5.6cm}}
\toprule
政策 & 实证命题（是什么） & 规范判断（应当怎样） \\
\midrule
最低工资 & 提高最低工资降低低技能工人的就业量（弹性可估计） & 应当/不应当实行最低工资（需福利标准与分配权重） \\
房产税 & 开征房产税降低投资性购房需求、改变地方财政收入结构 & 应当/不应当开征（效率与公平的权衡是价值问题） \\
消费券/补贴 & 消费券在短期内提高边际消费倾向高的群体的消费 & 应当/不应当发放（跨期效率、财政成本与目标人群是价值判断） \\
\bottomrule
\end{tabular}
\end{center}

注意左栏每一条原则上都可被数据\emph{证伪}；右栏每一条都必须明示其价值标准。两栏不可互相替代。

\begin{example}[最低工资：从实证估计到规范判断]
以美国青少年低技能劳动力市场为例。实证研究关心：最低工资提高 10\% 会导致就业下降多少？早期时间序列估计给出就业弹性约为 $-0.1$ 到 $-0.3$（提高工资显著减少就业）；而 Card \& Krueger（1994）利用新泽西与宾州快餐店的自然实验发现，新泽西最低工资从 \$4.25 提高到 \$5.05 后，就业变化并不显著为负，甚至略升。后续元分析（Neumark \& Wascher, 2007）指出：对青少年就业的弹性中位数约为 $-0.1$，且对低技能群体的负面影响更明显。

这些估计回答的是事实层与机制层。但“是否应当提高最低工资”仍取决于价值判断：若采用功利主义并赋予低收入者较高权重 $\alpha_i$，就业的轻微损失可能被低收入者工资收益抵消；若采用帕累托标准，则任何因失业受损的工人都构成反对理由，除非对他们进行补偿。实证结果缩小了规范争论的范围，却不能单独决定政策。
\end{example}

\subsection{实证与规范的互补}

实证与规范不是对手，而是完整研究的两个半部：
\[
\underbrace{\text{实证：预测政策后果}}_{\text{没有它，规范是空中楼阁}}
\;+\;
\underbrace{\text{规范：评价后果好坏}}_{\text{没有它，实证无政策含义}}
\;\Longrightarrow\; \text{完整的政策分析}.
\]
实证错误会让规范结论建立在沙上；规范模糊会让实证研究失去方向。博士生的写作规范是：每一段都先说清楚自己在哪一层。

\notehead{本节要点}
\begin{itemize}
  \item 实证：是什么/为什么/会怎样；三特征——概念清晰、机制明确、可检验；
  \item 规范：应当怎样；必须援引价值标准（帕累托、功利主义 $W=\sum\alpha_i u_i$、罗尔斯）；
  \item 权重 $\alpha_i$ 是价值立场而非事实；效率 $\neq$ 公平；
  \item 三层分离：事实层、机制层、价值层；越界是最常见错误。
\end{itemize}
